\subsection{Spinal Level Enumeration, Regions and Adjacency}
\label{sec:spn-1a}

This section defines the thoracolumbar intervertebral level vocabulary used by
the segmental spine model (issue \#12739, epic \#12735). It is pure
bookkeeping: no kinematics, no inertia and no OpenSim model are involved.

\paragraph{Scope.}
The module \texttt{src/shared/python/spine/levels.py} provides the enumeration
\texttt{SpinalLevel} (17 thoracolumbar disc levels, L5-S1 through T1-T2), the
region enumeration \texttt{SpinalRegion}, and the helpers
\texttt{levels\_between} and \texttt{parse\_level}. Cervical levels are out of
scope; the extension point is documented in the module docstring.

\paragraph{Ordering and Frames.}
Level $k$ with $k \in \{0,\dots,16\}$ denotes the disc space between an
inferior (caudal) vertebra $V_{\mathrm{inf}}$ and a superior (cranial)
vertebra $V_{\mathrm{sup}}$. The enum order is caudal to cranial:
\[
  \mathrm{L5\_S1} \;(k{=}0) \rightarrow \mathrm{L4\_L5} \rightarrow \cdots
  \rightarrow \mathrm{T1\_T2}\;(k{=}16).
\]
Adjacency is contiguous: for every $k < 16$,
$V_{\mathrm{sup}}(k) = V_{\mathrm{inf}}(k+1)$. The frame convention of the
transforms built on these levels (superior frame mapped into the inferior
frame, ISB axes) is defined in ADR \#12736 and is not restated here.

\paragraph{Functions.}
\begin{itemize}
  \item \texttt{inferior}, \texttt{superior}: vertebra labels, e.g.\
    $\mathrm{L4\_L5} \mapsto (\mathrm{L5}, \mathrm{L4})$ and
    $\mathrm{L5\_S1} \mapsto (\mathrm{S1}, \mathrm{L5})$.
  \item \texttt{region}: region of the superior vertebra, e.g.\
    $\mathrm{T12\_L1} \mapsto$ THORACIC and $\mathrm{L1\_L2} \mapsto$ LUMBAR.
  \item \texttt{ordinal}: zero-based caudal-to-cranial index.
  \item \texttt{levels\_between}$(c, r)$: the inclusive tuple
    $(\ell_{c}, \ell_{c+1}, \dots, \ell_{r})$ in caudal-to-cranial order.
    Precondition $\mathrm{ordinal}(c) \le \mathrm{ordinal}(r)$, otherwise
    \texttt{ValueError}; non-\texttt{SpinalLevel} arguments raise
    \texttt{TypeError}. Postcondition: first element $c$, last element $r$.
  \item \texttt{parse\_level}$(s)$: accepts \texttt{L4\_L5}, \texttt{L4-L5} and
    \texttt{L4/L5}, case-insensitive, with surrounding whitespace stripped.
    Non-\texttt{str} input raises \texttt{TypeError}; empty or unknown input
    raises \texttt{ValueError}.
\end{itemize}

\paragraph{Assumptions and Limitations.}
Labels follow the usual clinical convention for the thoracolumbar spine with
S1 as the caudal vertebra of L5-S1. The enumeration does not model sacral
segments beyond the S1 label, and \texttt{SpinalRegion.SACRAL} is defined but
no level uses it as its superior vertebra. Validation uses explicit
\texttt{TypeError}/\texttt{ValueError} checks so that the documented errors do
not depend on the DbC enforcement level.

\paragraph{Verification and Reproduction.}
Test evidence (30 unit tests, marker \texttt{unit}): contiguity of the adjacency
chain, the 17-level ordering, the five-member run \texttt{levels\_between}(L5\_S1,
L1\_L2), reversed-argument rejection, the three accepted spellings, and rejection
of malformed names and non-string input.
\begin{verbatim}
python3 -m pytest tests/shared/python/spine/test_levels.py -q
python3 -m ruff check src/shared/python/spine tests/shared/python/spine
python3 -m ruff format --check src/shared/python/spine tests/shared/python/spine
\end{verbatim}
